\section{Library Overview}
\label{sec:overview}

This section introduces the library through the user-facing C interface for a
guarded product of a double word and a binary64 value. 
first we present the guarded c function. 
then we present the double 

garbage:
We first define the double-word representation needed
to understand the operation.  We then present the C implementation of the
guarded interface, explain its nonreturning failure path and normal-return
path, and state what normal return guarantees: agreement with the functional
model, a proved accuracy bound, and a valid output representation.  Finally,
we present the numerical kernel and use its intermediate computations to
motivate the representation and range properties established by the check and
the treatment of gradual underflow in the numerical analysis.
Section~\ref{sec:specification} states the operation condition, error theorem,
guarded contract, and verification boundary precisely.

\subsection{The Guarded C Interface}
\label{sec:overview-interface}

Every guarded operation in the library follows the same control flow: an
operation-specific check is evaluated before the arithmetic kernel is called.
\Cref{fig:guarded-interface} shows the complete guarded function for
multiplication of a double word by a binary64 value.

\begin{figure}[t]
  \begin{lstlisting}[language=C]
void dw_times_fp_guarded(struct dword *z,
                        struct dword *x,
                        double y) {
    if (!dw_times_fp_check(x, y))
        exit(EXIT_FAILURE);

    dw_times_fp(z, x, y);
}
  \end{lstlisting}
  \caption{Guarded entry point for DW$\times$FP multiplication.  The
  arithmetic kernel is called only when the check succeeds.}
  \Description{A C function that exits when a runtime check fails and otherwise
  calls the double-word multiplication kernel.}
  \label{fig:guarded-interface}
\end{figure}

The argument \texttt{x} is a pointer to the structure storing the input double
word, \texttt{y} is the binary64 operand, and \texttt{z} is a pointer to the
structure in which the result is stored.  If the check fails, the function
calls \texttt{exit(EXIT\_FAILURE)}.  The arithmetic kernel is not called, the
function does not return normally, and no numerical result is returned.  If
the check succeeds, \texttt{dw\_mult\_fp} stores the computed pair in
\texttt{z} and the guarded function returns normally.  A normal return
therefore means that the check succeeded and the arithmetic kernel completed.

\begin{figure}[t]
  \begin{lstlisting}[language=C]
struct dword {
    double s;  /* high word */
    double t;  /* low word  */
};
  \end{lstlisting}
  \caption{C storage representation for a double word.}
  \Description{A C structure containing binary64 high-word and low-word
  fields.}
  \label{fig:dword-struct}
\end{figure}

\Cref{fig:dword-struct} shows the C structure used to store input and output
double words.  The fields \texttt{s} and \texttt{t} store the two binary64
components, referred to as the high and low words, respectively.

The structure specifies only how these components are stored.  The guarantee
on normal return concerns the value represented by the pair and the relation
between its components.

\subsection{Double-Word Representation}
\label{sec:overview-values}

A normally returning guarded call guarantees that a valid double-word
representation is stored in the output structure.  Let
$\operatorname{RN}(r)$ denote rounding a real number $r$ to binary64 using
round-to-nearest, ties-to-even.

\begin{definition}[Double word]
\label[definition]{def:double-word}
A double-word number $x$ is the unevaluated sum
\[
  x=x_h+x_\ell
\]
of two binary64 values $x_h$ and $x_\ell$ such that
\begin{equation}
  x_h=\operatorname{RN}(x_h+x_\ell).
  \label{eq:dw-representation}
\end{equation}
\end{definition}

A pair satisfying Equation~\eqref{eq:dw-representation} is a valid double
word.  Its components $x_h$ and $x_\ell$ are called the high and low words,
respectively.  The pair represents the real value $x=x_h+x_\ell$.  Since
$x_h=\operatorname{RN}(x)$, the low word is the residual
$x_\ell=x-x_h$ left after rounding $x$ to binary64.  Together, the two words
represent $x$ exactly.  The following example takes $x$ to be the exact
product of two binary64 values.

For example, let $a=b=1+2^{-52}$.  Their exact product is
\[
  ab=1+2^{-51}+2^{-104}.
\]
Rounding this product to binary64 gives
\[
  p_h=\operatorname{RN}(ab)=1+2^{-51},
\]
while an FMA recovers the discarded term:
\[
  p_\ell=\operatorname{FMA}(a,b,-p_h)=2^{-104}.
\]
Thus, $p_h+p_\ell=ab$.  The value $p_h$ is the high word, while $p_\ell$
retains the term discarded when the exact product is rounded to binary64.
The pair is a valid double word according to \Cref{def:double-word}.

Equation~\eqref{eq:dw-representation} relates the two components stored by a
normally returning guarded call but says nothing about the accuracy of the
operation.  Accuracy compares the value represented by the computed double
word with the exact real result.  Establishing an accuracy bound for the
running example, \texttt{dw\_mult\_fp\_guarded}, first requires specifying the
algorithm that computes the output pair.

\subsection{Accuracy of the Product of a Double Word and a Floating-Point Value}
\label{sec:overview-running}

\Cref{alg:dw-times-fp} presents \DWTimesFP{}, an algorithm for multiplying a
double word by a floating-point value \citep[Algorithm~9]{Joldes2017}.  The
inputs are a pair $(x_h,x_\ell)$, assumed to be a valid double-word
representation, and a floating-point value $y$.  The algorithm returns a pair
$(z_h,z_\ell)$, whose represented value $z_h+z_\ell$ approximates the exact
real product $(x_h+x_\ell)y$.

\begin{algorithm}[\DWTimesFP{}$(x_h,x_\ell,y)$]
\label[algorithm]{alg:dw-times-fp}
\begin{flushleft}
\(
  \begin{array}{@{}r@{\quad}l@{}}
    1: & (c_h,c_{\ell,1}) \leftarrow \FastTwoMult{}(x_h,y)\\
    2: & c_{\ell,3} \leftarrow \operatorname{FMA}(x_\ell,y,c_{\ell,1})\\
    3: & (z_h,z_\ell) \leftarrow \FastTwoSum{}(c_h,c_{\ell,3})\\
    4: & \mathbf{return}\ (z_h,z_\ell)
  \end{array}
\)
\end{flushleft}
\end{algorithm}

In \Cref{alg:dw-times-fp}, \FastTwoMult{} and \FastTwoSum{} are standard
building blocks for double-word arithmetic
\citep{Dekker1971,Joldes2017}.

The predicate $\DWTimesFPCond(x_h,x_\ell,y)$ denotes mathematical input and
range conditions sufficient to prove the following result.

\begin{theorem}[Accuracy and valid output of \DWTimesFP{}]
\label{thm:dw-times-fp-numerical}
Let $u=2^{-p}$ be unit roundoff for a radix-two format of precision $p\geq3$,
and let $\eta$ be its smallest positive subnormal value.  Let
\[
  (z_h,z_\ell)=\DWTimesFP{}(x_h,x_\ell,y).
\]
If $\DWTimesFPCond(x_h,x_\ell,y)$ holds, then
\begin{equation}
  \left|(z_h+z_\ell)-(x_h+x_\ell)y\right|
  \leq 2u^2\left|(x_h+x_\ell)y\right|+2.51\eta,
  \label{eq:dw-times-fp-mixed}
\end{equation}
and $(z_h,z_\ell)$ is a valid double word:
\[
  z_h=\operatorname{RN}(z_h+z_\ell).
\]
\end{theorem}

\citet{Joldes2017} proved a relative-error bound of $2u^2$ in a model
without underflow or overflow.  Our condition for \DWTimesFP{},
$\DWTimesFPCond(x_h,x_\ell,y)$, prevents overflow but permits intermediate
values near the underflow threshold.  For such values, the usual
relative-error model does not apply uniformly, so we
bound rounding error in absolute terms using the smallest positive subnormal
value, $\eta$.  For binary64, $\eta=2^{-1074}$.  Propagating these absolute
bounds through the algorithm gives the absolute error term $2.51\eta$.

% On Line~1, \FastTwoMult{}$(x_h,y)$ returns $(c_h,c_{\ell,1})$, where $c_h$
% is the computed product and $c_{\ell,1}$ is a correction for its rounding
% error.  When the exact correction is representable, the two terms recover the
% exact leading product:
%\[
%  c_h+c_{\ell,1}=x_h\cdot y.
%\]
%Near underflow, this equality may fail; the absolute term in
%\Cref{eq:dw-times-fp-mixed} accounts for the resulting error.

%Line~2 uses a fused multiply-add to combine the residual $c_{\ell,1}$ with
%the contribution from the input low word, producing $c_{\ell,3}$.  However,
%$c_{\ell,3}$ is only a candidate low word: it combines the rounding error
%from the high-word product with the contribution $x_\ell\cdot y$ from the
%input low word.  This combined term may be large enough that
%\[
%  \operatorname{RN}(c_h+c_{\ell,3})\neq c_h.
%\]
%The pair $(c_h,c_{\ell,3})$ may therefore fail to satisfy the double-word
%relation in \Cref{def:double-word}.

%Line~3 therefore applies \FastTwoSum{} to $(c_h,c_{\ell,3})$ to construct a
%valid double word, returning $(z_h,z_\ell)$ such that
%\[
%  z_h=\operatorname{RN}(c_h+c_{\ell,3})
%  \qquad\text{and}\qquad
%  z_h+z_\ell=c_h+c_{\ell,3}.
%\]
%Consequently,
%\[
%  z_h=\operatorname{RN}(z_h+z_\ell),
%\]
%so the returned pair is a valid double word according to
%\Cref{def:double-word}.

To use \Cref{thm:dw-times-fp-numerical} at the C interface, a successful
executable check must establish $\DWTimesFPCond$.  The guard in
\Cref{fig:guarded-interface} runs this check before invoking the arithmetic
kernel.

\subsection{The Runtime Check}
\label{sec:overview-guard}

\Cref{fig:dw-times-fp-check} gives the implementation of this check.

\begin{figure}[t]
  \begin{lstlisting}[language=C]
int dw_times_fp_check(struct dword *x,
                     double y) {
    int fin_xh = isfinite_double(x->s);
    int fin_xl = isfinite_double(x->t);
    int fin_y  = isfinite_double(y);
    int dwx = fin_xh && fin_xl &&
              (x->s + x->t == x->s);

    double z = x->s * y;
    double q_prod = ldexp(1.0, 1023);
    int prod_bnd = fabs(z) < q_prod;

    return dwx && fin_y && prod_bnd;
}
  \end{lstlisting}
  \caption{Runtime check for the product of a double word and a binary64
  value.}
  \Description{C code that checks input finiteness, the double-word relation,
  and an upper bound on the rounded leading product.}
  \label{fig:dw-times-fp-check}
\end{figure}

When \texttt{x} stores $(x_h,x_\ell)$, we prove
\[
  \texttt{dw\_mult\_fp\_check}(x,y)=1
  \quad\Longrightarrow\quad
  \DWTimesFPCond(x_h,x_\ell,y).
\]
Thus, a successful check establishes the sufficient condition used in
\Cref{thm:dw-times-fp-numerical}.

Lines~3--7 establish the required properties of the inputs.  The finiteness
tests exclude infinities and NaNs.  Line~7 of
\Cref{fig:dw-times-fp-check} checks the relation between the components of the
input pair by evaluating
\begin{lstlisting}[language=C,numbers=none,
                   backgroundcolor=\color{gray!5}]
x->s + x->t == x->s
\end{lstlisting}
Under the strict IEEE~754 semantics required by the library, the addition
computes $\operatorname{RN}(x_h+x_\ell)$.  A true comparison therefore
establishes the double-word relation in \Cref{def:double-word}.

Lines~9--11 establish the upper-range bound.  The code computes
$z=\operatorname{RN}(x_hy)$ and checks that $|z|<2^{1023}$.  The corresponding
condition in $\DWTimesFPCond$ concerns the exact product $x_hy$, which cannot
be computed directly using binary64 arithmetic.  However, $2^{1023}$ is
exactly representable in binary64.  By monotonicity of round-to-nearest, an
exact product with magnitude at least $2^{1023}$ cannot round to a value with
smaller magnitude.  Passing the strict test therefore establishes
$|x_hy|<2^{1023}$, providing the upper-range bound used to prevent overflow
in \DWTimesFP{}.

The implication is intentionally one way: the strict comparison may
conservatively reject an input for which $\DWTimesFPCond$ holds.  At the lower
end of the format, the check imposes no magnitude restriction.  The absolute
term in \Cref{eq:dw-times-fp-mixed} accounts for rounding in this regime.

On the successful branch of \Cref{fig:guarded-interface}, soundness of the check
supplies the hypothesis of the numerical theorem before the arithmetic
kernel is invoked.  Completing the guarantee for a normally returning call
then requires a refinement proof showing that the C kernel stores the pair
returned by \DWTimesFP{}.

\subsection{Establishing the Guarded Contract}
\label{sec:overview-pipeline}

We can now state the guarantee provided by the running example at the C
interface.  Suppose that \texttt{x} initially stores $(x_h,x_\ell)$ and that
\texttt{dw\_mult\_fp\_guarded} returns normally with \texttt{z} storing
$(z_h,z_\ell)$.  The stored pair is the binary64 pair specified by
\DWTimesFP{} on the same inputs, its represented value satisfies the bound in
\Cref{eq:dw-times-fp-mixed}, and it is a valid double word according to
\Cref{def:double-word}.  Callers do not establish
$\DWTimesFPCond(x_h,x_\ell,y)$; the guarded interface establishes a sufficient
condition before invoking the arithmetic kernel.

\definecolor{flocqshade}{RGB}{222,236,249}
\definecolor{vstshade}{RGB}{224,241,229}
\definecolor{compcertshade}{RGB}{252,232,211}
\definecolor{artifactshade}{RGB}{243,243,243}
\newcommand{\diagramboxtitle}[1]{%
  \makebox[\linewidth][c]{{\normalsize\textbf{#1}}}}

\begin{figure*}[t]
\centering
\begin{tikzpicture}[
  x=1cm,
  y=1cm,
  node distance=4mm,
  flow/.style={-{Latex[length=2.1mm]},thick},
  connector/.style={thick},
  artifact/.style={draw,rounded corners=2pt,fill=artifactshade,
    align=left,inner sep=5pt,font=\small},
  flocqproof/.style={draw,rounded corners=2pt,fill=flocqshade,
    align=left,inner sep=5pt,font=\small},
  vstproof/.style={draw,rounded corners=2pt,fill=vstshade,
    align=left,inner sep=5pt,font=\small},
  compcertproof/.style={draw,rounded corners=2pt,fill=compcertshade,
    align=left,inner sep=5pt,font=\small},
  arrowlabel/.style={font=\scriptsize,fill=white,inner sep=1.5pt,align=center}
]

\node[artifact,text width=7.6cm,anchor=north] (csource) at (-4.15,0)
  {\diagramboxtitle{Handwritten C sources}\\[2pt]
   \makebox[\linewidth][c]{arithmetic kernel, runtime check, and guarded interface}};

\node[artifact,text width=7.6cm,anchor=north] (model) at (4.15,0)
  {\diagramboxtitle{Rocq models and definitions}\\[2pt]
   \makebox[\linewidth][c]{\DWArithOp{} and $\DWArithOpCond$}};

\node[vstproof,text width=5.55cm,minimum height=2cm,
      anchor=north] (checkproof) at (-5.9,-2.1)
  {\diagramboxtitle{Soundness of the check}\\[3pt]
   \makebox[\linewidth][c]{successful check $\Rightarrow \DWArithOpCond$}};

\node[vstproof,text width=5.2cm,minimum height=2cm,
      anchor=north] (kernelproof) at (0,-2.1)
  {\diagramboxtitle{C refinement proof}\\[3pt]
   \makebox[\linewidth][c]{C kernel satisfies \DWArithOp{} spec}};

\node[flocqproof,text width=5.55cm,minimum height=2cm,
      anchor=north] (numproof) at (5.9,-2.1)
  {\diagramboxtitle{Numerical guarantees}\\[3pt]
   $\DWArithOpCond \Rightarrow$ result is accurate\\[-1pt]
   and a valid double-word};

\coordinate (cfork) at (-4.15,-1.75);
\coordinate (mfork) at (4.15,-1.75);
\draw[connector] (csource.south) -- (cfork);
\draw[flow] (cfork) -| (checkproof.north);
\draw[flow] (cfork) -| ([xshift=-11mm]kernelproof.north);
\draw[connector] (model.south) -- (mfork);
\draw[flow] (mfork) -| ([xshift=11mm]kernelproof.north);
\draw[flow] (mfork) -| (numproof.north);

\node[vstproof,text width=14cm,anchor=north] (wrapper)
      at (0,-4.7)
  {\diagramboxtitle{Guarded interface proof}\\[3pt]
   \makebox[\linewidth][c]{successful branch combines the three results;
   failure branch exits before calling the kernel}};
\draw[flow] (checkproof.south) -- ([xshift=-59mm]wrapper.north);
\draw[flow] (kernelproof.south) -- (wrapper.north);
\draw[flow] (numproof.south) -- ([xshift=59mm]wrapper.north);

\node[artifact,text width=10.5cm,below=of wrapper] (cguarantee)
  {\diagramboxtitle{Verified C interface}\\[2pt]
   check failure $\Rightarrow$ nonreturning error path\\[-1pt]
   normal return $\Rightarrow$ result refines functional-model, satisfies
   accuracy guarantee, and is a valid double-double};
\draw[flow] (wrapper) -- (cguarantee);

\node[compcertproof,text width=6.2cm,below=8mm of cguarantee] (compcert)
  {\diagramboxtitle{CompCert compiler correctness}};

\node[artifact,text width=5.8cm,below=of compcert] (assembly)
  {\diagramboxtitle{Generated assembly}\\[-1pt]
   \makebox[\linewidth][c]{preserves the verified C behavior}};
\draw[flow] (compcert) -- (assembly);

\node[draw,dashed,rounded corners=2pt,
      fit=(compcert)(assembly),inner xsep=6mm,inner ysep=4mm]
      (compcertgroup) {};
\node[font=\small,fill=white,inner xsep=3pt]
      at (compcertgroup.north) {When compiled with CompCert};
\draw[flow] ([xshift=-30mm]cguarantee.south) --
            ([xshift=-30mm]compcertgroup.north);
\draw[flow] ([xshift=30mm]cguarantee.south) --
            ([xshift=30mm]compcertgroup.north);

\end{tikzpicture}
\caption{For each guarded operation, soundness of the check, the C refinement
proof, and the numerical theorem combine to establish the guarded C
contract.  When CompCert is used, compiler correctness extends this guarantee
to generated assembly.  Green boxes show VST proofs, the blue box shows the
Flocq proof, and the orange box shows the CompCert correctness theorem.}
\Description{A top-to-bottom diagram.  C sources and Rocq definitions lead
to three verified results: soundness of the runtime check, a refinement proof
showing that the C kernel implements its functional model, and an
operation-specific numerical theorem proved in Rocq using Flocq.  A VST proof
combines these results into the guarded C contract.  CompCert correctness
preserves that behavior in generated assembly when CompCert is used.}
\label{fig:overview-architecture}
\end{figure*}

Three separately proved results establish this contract.  Soundness of the
check shows that a successful call to \texttt{dw\_mult\_fp\_check} establishes
$\DWTimesFPCond(x_h,x_\ell,y)$.  A VST refinement proof for
\texttt{dw\_mult\_fp} shows that the C kernel stores the pair returned by
\DWTimesFP{} on the same inputs.  Finally,
\Cref{thm:dw-times-fp-numerical} gives the error bound and valid double-word
result whenever $\DWTimesFPCond$ holds.

The VST proof of the guarded function combines these results by following the
two branches in \Cref{fig:guarded-interface}.  A failed check reaches the
nonreturning error path without invoking the kernel.  On the successful
branch, soundness of the check supplies the condition needed by the numerical
theorem, while the refinement proof identifies the stored result with the output
of \DWTimesFP{}.  \Cref{fig:overview-architecture} generalizes this composition
across the library.  In the figure, \DWArithOp{} denotes the functional model
for an arbitrary operation and $\DWArithOpCond$ denotes its sufficient
condition.

\Cref{tab:operation-preview} summarizes the operations and normal-return
accuracy guarantees provided by the library.  Let $r$ denote the exact real
result and let $\widehat r=z_h+z_\ell$ denote the value represented by the
returned pair.  Every normally returning call also guarantees that the pair
is the result of the corresponding functional model and is a valid double
word.

\begin{table}[t]
\centering
\caption{Operations and normal-return accuracy guarantees provided by the
library.  For each operation, $r$ denotes the exact real result and
$\widehat r=z_h+z_\ell$ denotes the real value represented by the returned
pair.  The value $u=2^{-53}$ is the unit roundoff and $\eta=2^{-1074}$ is the
smallest positive subnormal value for binary64.  Normal return also guarantees
that the pair is a valid double-double.}
\Description{The library operations, their exact real results, and the
accuracy guarantees provided when their guarded interfaces return normally.}
\label{tab:operation-preview}
\small
\begin{tabular}{@{}lll@{}}
\toprule
Operation & Exact result &
  \shortstack[l]{Normal-return\\accuracy guarantee} \\
\midrule
\TwoSum{}         & $a+b$ & $\widehat r=r$ (exact) \\
\FastTwoSum{}     & $a+b$ & $\widehat r=r$ (exact) \\
\FastTwoMult{}    & $ab$  & $\widehat r=r$ (exact) \\
\addlinespace
\textsc{DWPlusFP} & $x+y$ &
  $|\widehat r-r|\leq 2u^2|r|$ \\
\textsc{DWPlusDW} & $x+y$ &
  $|\widehat r-r|\leq \dfrac{3u^2}{1-4u}|r|$ \\
\DWTimesFP{}       & $xy$ &
  $|\widehat r-r|\leq 2u^2|r|+2.51\eta$ \\
\textsc{DWTimesDW} & $xy$ &
  $|\widehat r-r|\leq 4u^2|r|+2.5\eta$ \\
\textsc{DWDivFP}  & $x/y$ &
  $|\widehat r-r|\leq 3u^2|r|+\eta$ \\
\textsc{DWDivDW}  & $x/y$ &
  $|\widehat r-r|\leq 12.33u^2|r|+2\eta$ \\
\bottomrule
\end{tabular}
\end{table}

In this paper, \emph{end-to-end} means that machine-checked proofs connect the
numerical guarantees for the functional model to the contract exposed by the
linked guarded C interface and, when CompCert is used, to the behavior of the
generated assembly.  The numerical and C proofs require strict IEEE~754
semantics.  Any compiler used to build the library must preserve these
semantics.  CompCert satisfies this requirement, and its correctness theorem
additionally preserves the verified C behavior during compilation.  Code
produced by another compiler is not covered by this compilation theorem.
\Cref{sec:specification} states the mathematical condition, C contract, and
verification boundary precisely.
